\section{The reassembly certificate lattice}
\label{sec:tpc62-certificate-lattice}

The same physical reassembly map admits several lower certificates.  They
do not all contain the same information.  This section places the exact
shorted constant, exact missing-to-represented ratios, the restricted
native Gram, a two-stage residual certificate, and the ambient missing
degree in one partial order.  The order is valid only when the scope---all
coefficient vectors or one distinguished vector---is held fixed.

\subsection{The exact shorted constant}

Let \(\cK,\cH\) be finite-dimensional complex Hilbert spaces and let
\begin{equation}
 A,M:\cK\longrightarrow\cH,
 \qquad B:=A+M,
 \qquad A\ne0.
 \label{eq:tpc62-lattice-AMB}
\end{equation}
The represented map is \(A\), the missing signed contribution is \(M\),
and \(B\) is the physical reassembly.  Its optimal coefficient-uniform
constant relative to the represented norm is
\begin{equation}
 \delta_{\rm sh}(A,M)
 :=\inf_{Az\ne0}\frac{\|Bz\|^2}{\|Az\|^2}.
 \label{eq:tpc62-lattice-delta-definition}
\end{equation}
In the notation of the master gate, this is exactly
\(\alpha_{A^*A}(A+M)\).  We repeat the specialization here so that every
node of the certificate lattice is defined in one place.

Put
\begin{equation}
 N:=\Ker A,
 \qquad S:=N^\perp,
 \qquad A_S:=A|_S,
 \label{eq:tpc62-lattice-NS}
\end{equation}
and write
\begin{equation}
 B_S:=B|_S,
 \qquad B_N:=B|_N=M|_N,
 \qquad P:=P_{(\Ran B_N)^\perp}.
 \label{eq:tpc62-lattice-BSBP}
\end{equation}
Finally set
\begin{equation}
 R_S:=A_S^*A_S>0,
 \qquad Q_S:=B_S^*PB_S\ge0.
 \label{eq:tpc62-lattice-RQ}
\end{equation}

\begin{theorem}[Exact shorted reassembly constant]
\label{thm:tpc62-lattice-exact-shorted}
For the maps in \eqref{eq:tpc62-lattice-AMB},
\begin{equation}
 \boxed{
 \delta_{\rm sh}(A,M)
 =\lambda_{\min}\!\left(
   R_S^{-1/2}Q_SR_S^{-1/2}
  \right).}
 \label{eq:tpc62-lattice-exact-shorted-formula}
\end{equation}
This is an identity, not a sufficient estimate.
\end{theorem}

\begin{proof}
Every \(z\in\cK\) has a unique decomposition \(z=s+n\), with
\(s\in S\) and \(n\in N\).  The denominator depends only on \(s\):
\(Az=A_Ss\).  For this fixed \(s\), least squares gives
\begin{align*}
 \inf_{n\in N}\|B_Ss+B_Nn\|^2
 &=\dist(B_Ss,\Ran B_N)^2\\
 &=\|PB_Ss\|^2
  =\langle Q_Ss,s\rangle.
\end{align*}
Since \(A_S\) is injective, \(R_S\) is positive definite.  Taking the
infimum over \(s\ne0\) and applying Rayleigh--Ritz proves
\eqref{eq:tpc62-lattice-exact-shorted-formula}.
\end{proof}

The projection \(P\) is essential: directions invisible to \(A\) may
still be visible to \(B\), and can then cancel part of \(B_Ss\) without
changing the denominator.  Any lower certificate which ignores this
nuisance minimization is not the exact coefficient-uniform constant.

\subsection{Exact norm ratios and the triangle certificate}

For \(Az\ne0\), define the pointwise missing ratio and physical ratio
\begin{equation}
 q(z):=\frac{\|Mz\|^2}{\|Az\|^2},
 \qquad
 r(z):=\frac{\|(A+M)z\|^2}{\|Az\|^2}.
 \label{eq:tpc62-pointwise-q-r}
\end{equation}
The corresponding uniform missing ratio is the extended-real number
\begin{equation}
 q_{\rm unif}:=\sup_{Az\ne0}q(z)\in[0,\infty].
 \label{eq:tpc62-uniform-q}
\end{equation}
In finite dimension there is an exact dichotomy:
\begin{equation}
 \boxed{
 q_{\rm unif}<\infty
 \quad\Longleftrightarrow\quad
 \Ker A\subseteq\Ker M.}
 \label{eq:tpc62-q-finite-kernel-test}
\end{equation}
Indeed, if \(n\in\Ker A\) but \(Mn\ne0\), choose \(s\) with \(As\ne0\).
Then \(A(s+tn)=As\), whereas
\(\|M(s+tn)\|\to\infty\) as \(|t|\to\infty\), so
\(q_{\rm unif}=\infty\).  Conversely, if \(M\) vanishes on
\(\Ker A\), both numerator and denominator depend only on
\(s\in(\Ker A)^\perp\); compactness of the corresponding normalized
sphere gives a finite supremum.  The exact shorted constant remains well
defined even in the infinite-\(q_{\rm unif}\) case.

Let
\begin{equation}
 \varphi(t):=(1-\sqrt t)_+^2
 \quad(0\le t<\infty),
 \qquad
 \varphi(\infty):=0.
 \label{eq:tpc62-triangle-profile}
\end{equation}
This function is nonincreasing.

\begin{proposition}[Pointwise and uniform triangle bounds]
\label{prop:tpc62-q-certificates}
For every \(z\) with \(Az\ne0\),
\begin{equation}
 \boxed{r(z)\ge\varphi(q(z)).}
 \label{eq:tpc62-pointwise-q-certificate}
\end{equation}
Uniformly over all coefficient directions,
\begin{equation}
 \boxed{
 \delta_{\rm sh}(A,M)\ge\varphi(q_{\rm unif}).}
 \label{eq:tpc62-uniform-q-certificate}
\end{equation}
Neither inequality is an identity in general.
\end{proposition}

\begin{proof}
The reverse triangle inequality gives
\[
 \|(A+M)z\|
 \ge\bigl(\|Az\|-\|Mz\|\bigr)_+.
\]
Divide by \(\|Az\|\), square, and use
\eqref{eq:tpc62-pointwise-q-r} to obtain
\eqref{eq:tpc62-pointwise-q-certificate}.  Since
\(q(z)\le q_{\rm unif}\) and \(\varphi\) is nonincreasing, taking the
infimum over \(z\) proves
\eqref{eq:tpc62-uniform-q-certificate}.
\end{proof}

Two elementary examples show why the shorted constant is strictly more
informative.

\begin{example}[Angle information is discarded]
\label{ex:tpc62-q-angle-loss}
Take \(\cK=\cH=\C\), \(Az=z\), and \(Mz=iz\).  Then
\[
 q_{\rm unif}=1,
 \qquad
 \varphi(q_{\rm unif})=0,
 \qquad
 \delta_{\rm sh}(A,M)=|1+i|^2=2.
\]
The norm of \(Mz\) alone does not record its favorable angle with \(Az\).
\end{example}

\begin{example}[An infinite uniform \(q\) can coexist with positive
shorted reassembly]
\label{ex:tpc62-q-infinite-shorted-positive}
Take \(\cK=\cH=\C^2\) and
\[
 A(x,n)=(x,0),
 \qquad
 M(x,n)=(0,n).
\]
Then \(q_{\rm unif}=\infty\), so the triangle certificate is zero, but
\[
 \frac{\|(A+M)(x,n)\|^2}{\|A(x,n)\|^2}
 =1+\frac{|n|^2}{|x|^2}\ge1.
\]
Thus \(\delta_{\rm sh}(A,M)=1\).  The shorting projection removes the
orthogonal output generated by the \(A\)-invisible coordinate.
\end{example}

\subsection{The restricted native constant}

Let \(\mathbf D\ge0\) be the raw atomic diagonal on \(\cK\), and write
\begin{equation}
 \cD(z):=\langle\mathbf Dz,z\rangle.
 \label{eq:tpc62-raw-atomic-form}
\end{equation}
We assume
\begin{equation}
 \Ker\mathbf D\subseteq\Ker M,
 \label{eq:tpc62-D-kernel-compatible}
\end{equation}
as holds when \(\mathbf D=L^*L\) is formed from the literal active atomic
lift and \(M\) is obtained by grouping those atoms.  Put
\(S_D:=\Ran\mathbf D=(\Ker\mathbf D)^\perp\).

If \(S_D\ne\{0\}\), define
\begin{equation}
 \boxed{
 \kappa_M
 :=\lambda_{\max}\!\left(
   \left.
   \mathbf D^{\dagger/2}M^*M\mathbf D^{\dagger/2}
   \right|_{S_D}
  \right).}
 \label{eq:tpc62-restricted-kappa-definition}
\end{equation}
If \(S_D=\{0\}\), set \(\kappa_M=0\).  In the Rayleigh formula below the
supremum over the empty set is likewise defined to be zero.
Rayleigh--Ritz gives the exact equivalent formulas
\begin{align}
 \kappa_M
 &=\sup_{\cD(z)>0}\frac{\|Mz\|^2}{\cD(z)},
 \label{eq:tpc62-restricted-kappa-Rayleigh}\\
 &=\inf\{\kappa\ge0:M^*M\preceq\kappa\mathbf D\}.
 \label{eq:tpc62-restricted-kappa-Loewner}
\end{align}
Thus \(\kappa_M\) is the least coefficient-uniform constant in
\begin{equation}
 \|Mz\|^2\le\kappa_M\cD(z).
 \label{eq:tpc62-restricted-kappa-bound}
\end{equation}

Indeed, for \(\cD(z)>0\), set \(x=\mathbf D^{1/2}z\in S_D\).  By
\eqref{eq:tpc62-D-kernel-compatible}, replacing \(z\) by
\(\mathbf D^{\dagger/2}x\) does not change \(Mz\), and hence
\[
 \frac{\|Mz\|^2}{\cD(z)}
 =\frac{\langle
  \mathbf D^{\dagger/2}M^*M\mathbf D^{\dagger/2}x,x
  \rangle}{\|x\|^2}.
\]
Taking the maximum Rayleigh quotient proves
\eqref{eq:tpc62-restricted-kappa-Rayleigh}.  A quadratic-form bound for
all \(z\) is equivalent to
\(M^*M\preceq\kappa\mathbf D\), which proves
\eqref{eq:tpc62-restricted-kappa-Loewner} and the asserted optimality.

To compare \(\cD\) with represented mass, put
\begin{equation}
 \eta(z):=\frac{\cD(z)}{\|Az\|^2}
 \quad(Az\ne0),
 \qquad
 \eta_{\rm unif}:=\sup_{Az\ne0}\eta(z)\in[0,\infty].
 \label{eq:tpc62-eta-pointwise-uniform}
\end{equation}
The exact pointwise restricted quotient is
\begin{equation}
 \kappa_M(z):=
 \begin{cases}
 \|Mz\|^2/\cD(z),&\cD(z)>0,\\
 0,&\cD(z)=0.
 \end{cases}
 \label{eq:tpc62-pointwise-kappa-definition}
\end{equation}
The second line is consistent with
\eqref{eq:tpc62-D-kernel-compatible}.  Therefore
\begin{equation}
 q(z)=\kappa_M(z)\eta(z),
 \qquad
 0\le\kappa_M(z)\le\kappa_M.
 \label{eq:tpc62-pointwise-factorization}
\end{equation}

\subsection{Residual and ambient upper bounds}

There are two standard ways to upper-bound \(\kappa_M\), and they need
not be comparable with one another.

First suppose the missing map factors through a residual-labelled space,
\begin{equation}
 M=\mathsf E T,
 \qquad
 T:\cK\to\cR,
 \qquad
 \mathsf E:\cR\to\cH
 \label{eq:tpc62-residual-factorization}
\end{equation}
and assume \(\Ker\mathbf D\subseteq\Ker T\), as for a residual grouping
of the same literal atomic lift.  Define
\begin{align}
 \kappa_{\rm res}
 &:=\sup_{\cD(z)>0}\frac{\|Tz\|^2}{\cD(z)},
 \label{eq:tpc62-kappa-res-definition}\\
 \beta_M
 &:=\sup_{0\ne y\in\Ran T}
       \frac{\|\mathsf Ey\|^2}{\|y\|^2}.
 \label{eq:tpc62-beta-definition}
\end{align}
Zero maps are assigned constant zero.  Direct substitution gives
\begin{equation}
 \boxed{\kappa_M\le\kappa_{\rm res}\beta_M.}
 \label{eq:tpc62-residual-product-bound}
\end{equation}
The restriction of \(\beta_M\) to the actual image \(\Ran T\) is
essential; adjoining unused residual directions can only make the
certificate coarser.

Second, suppose the literal atomic factorization is
\begin{equation}
 \mathbf D=L^*L,
 \qquad
 M=\mathsf C L,
 \label{eq:tpc62-ambient-atomic-factorization}
\end{equation}
where \(\mathsf C\) sums the active atomic coordinates in each native
output fiber.  If \(\cF_i\) is the active fiber at native output \(i\),
set
\begin{equation}
 d_{\rm missing}:=\max_i|\cF_i|.
 \label{eq:tpc62-ambient-missing-degree}
\end{equation}
When there is no active fiber, this maximum is defined to be zero.
Fiberwise Cauchy--Schwarz is sharp on the free atomic space and gives
\(\|\mathsf C\|^2=d_{\rm missing}\).  Since
\(L\mathbf D^{\dagger/2}|_{S_D}\) is an isometry,
\begin{equation}
 \boxed{\kappa_M\le d_{\rm missing}.}
 \label{eq:tpc62-ambient-degree-bound}
\end{equation}
Both \eqref{eq:tpc62-residual-product-bound} and
\eqref{eq:tpc62-ambient-degree-bound} are upper certificates for the same
optimal restricted constant; neither is a second loss to be added to it.

\subsection{The dominance partial order}

For lower certificates in the same scope, say that \(C_1\) dominates
\(C_2\) when \(C_1\ge C_2\) for every admissible instance.  Thus the
larger number is the stronger guaranteed lower bound.

\begin{theorem}[Uniform certificate lattice]
\label{thm:tpc62-uniform-certificate-lattice}
Assume \(\eta_{\rm unif}<\infty\).  Then
\begin{equation}
 q_{\rm unif}
 \le\kappa_M\eta_{\rm unif}
 \le
 \begin{cases}
  \kappa_{\rm res}\beta_M\eta_{\rm unif},\\
  d_{\rm missing}\eta_{\rm unif}.
 \end{cases}
 \label{eq:tpc62-upper-constant-lattice}
\end{equation}
Consequently the corresponding coefficient-uniform lower certificates
satisfy
\begin{equation}
 \boxed{
 \delta_{\rm sh}(A,M)
 \quad\ge\quad \varphi(q_{\rm unif})
 \quad\ge\quad \varphi(\kappa_M\eta_{\rm unif})
 \quad\ge\quad
 \begin{cases}
  \varphi(\kappa_{\rm res}\beta_M\eta_{\rm unif}),\\
  \varphi(d_{\rm missing}\eta_{\rm unif}).
 \end{cases}}
 \label{eq:tpc62-lower-certificate-lattice}
\end{equation}
The restricted certificate therefore dominates each coarse certificate.
The residual-product and ambient-degree branches are generally
incomparable; they form a fork, not a linear chain.
\end{theorem}

\begin{proof}
By \eqref{eq:tpc62-restricted-kappa-bound}, for every \(Az\ne0\),
\[
 q(z)
 =\frac{\|Mz\|^2}{\|Az\|^2}
 \le\kappa_M\frac{\cD(z)}{\|Az\|^2}
 \le\kappa_M\eta_{\rm unif}.
\]
Taking the supremum proves the first inequality in
\eqref{eq:tpc62-upper-constant-lattice}.  The remaining two follow from
\eqref{eq:tpc62-residual-product-bound} and
\eqref{eq:tpc62-ambient-degree-bound}.  Apply the nonincreasing function
\(\varphi\), and use
\cref{prop:tpc62-q-certificates}, to obtain
\eqref{eq:tpc62-lower-certificate-lattice}.  The explicit strict examples
below prove incomparability of the two terminal branches.
\end{proof}

The lattice records three distinct information losses.  Passing from
\(\delta_{\rm sh}\) to \(q_{\rm unif}\) discards the angle and the exact
nuisance elimination.  Passing from \(q_{\rm unif}\) to
\(\kappa_M\eta_{\rm unif}\) separates two suprema which may be attained
on different coefficient directions.  Passing from \(\kappa_M\) to
either coarse branch enlarges the actual coefficient image.

\subsection{Explicit incomparability models}

We now construct two literal finite grouping models.  In both, the raw
atomic norm is \(\cD(z)=\|Lz\|^2\), residual labels are retained by \(T\),
and \(\mathsf E\) erases them.  Hence the examples compare the actual
constants in \eqref{eq:tpc62-residual-product-bound} and
\eqref{eq:tpc62-ambient-degree-bound}, rather than arbitrary numerical
upper bounds.

\begin{example}[The residual product can be strictly smaller than the
ambient degree]
\label{ex:tpc62-residual-better-than-ambient}
Let \(\cK=\C\), let the free atomic space be \(\C^2\), and let two atoms
share one native output but carry distinct residual labels.  Define
\begin{equation}
 Lz=\frac{z}{\sqrt2}(1,i),
 \qquad
 T=L,
 \qquad
 \mathsf E(y_1,y_2)=y_1+y_2.
 \label{eq:tpc62-residual-better-model}
\end{equation}
Then \(\mathbf D=L^*L=1\), \(\kappa_{\rm res}=1\), and
\(\Ran T=\C(1,i)\).  Therefore
\begin{equation}
 \beta_M
 =\frac{|1+i|^2}{2}=1,
 \qquad
 \kappa_{\rm res}\beta_M=1.
 \label{eq:tpc62-residual-better-product}
\end{equation}
The free native grouping has a fiber of size two, so
\(d_{\rm missing}=2\).  The physical map is
\(Mz=(1+i)z/\sqrt2\), and hence \(\kappa_M=1\).  Thus
\begin{equation}
 \kappa_M=\kappa_{\rm res}\beta_M=1<2=d_{\rm missing}.
 \label{eq:tpc62-residual-strictly-better}
\end{equation}
Both atoms are active; the saving comes from restricting the erasure to
the actual one-dimensional coefficient image, not from deleting an atom.
\end{example}

\begin{example}[The ambient degree can be strictly smaller than the
residual product]
\label{ex:tpc62-ambient-better-than-residual}
Fix integers \(r,k\ge2\).  Take a free atomic space with two native
fibers.  Fiber \(i=1\) contains \(r\) atoms, all with residual label
\(s=1\).  Fiber \(i=2\) contains \(k\) atoms with distinct residual
labels \(s=1,\ldots,k\).  Let \(\cK\) be the full atomic space and
\(L=I\), so \(\mathbf D=I\).

Let \(T\) first sum atoms having the same pair \((i,s)\), and let
\(\mathsf E\) then sum residual labels at each fixed \(i\).  The first
map has squared norm
\begin{equation}
 \kappa_{\rm res}=r,
 \label{eq:tpc62-residual-first-stage-r}
\end{equation}
attained in the size-\(r\) group at \((1,1)\).  The map \(T\) is onto its
\((1+k)\)-dimensional residual output: the first grouped coordinate and
each of the \(k\) singleton coordinates can be prescribed independently.
Consequently
\begin{equation}
 \beta_M=k,
 \qquad
 \kappa_{\rm res}\beta_M=rk.
 \label{eq:tpc62-residual-product-rk}
\end{equation}
The composed map \(M=\mathsf ET\) is simply native grouping in the two
fibers.  Its squared norm and its ambient degree are
\begin{equation}
 \kappa_M=d_{\rm missing}=\max\{r,k\}.
 \label{eq:tpc62-ambient-max-rk}
\end{equation}
Thus
\begin{equation}
 d_{\rm missing}=\max\{r,k\}<rk
 =\kappa_{\rm res}\beta_M.
 \label{eq:tpc62-ambient-strictly-better}
\end{equation}
The two stage norms are attained in different native fibers, which is why
multiplying them loses information.
\end{example}

Together,
\cref{ex:tpc62-residual-better-than-ambient,ex:tpc62-ambient-better-than-residual}
prove that no universal arrow exists between
\(\kappa_{\rm res}\beta_M\) and \(d_{\rm missing}\).  In contrast, the
arrows from both of them down to the optimal restricted constant
\(\kappa_M\) are universal and sharp.

\subsection{Distinguished vectors and scope preservation}

The preceding lattice is coefficient-uniform.  For one distinguished
physical vector \(z_*\) with \(Az_*\ne0\), the correct pointwise lattice
is different.

\begin{proposition}[Distinguished-vector certificate lattice]
\label{prop:tpc62-distinguished-certificate-lattice}
Let \(z_*\) satisfy \(Az_*\ne0\).  Then
\begin{align}
 r(z_*)
 &\ge\varphi(q(z_*))
  =\varphi\!\left(\kappa_M(z_*)\eta(z_*)\right)
 \label{eq:tpc62-distinguished-exact-q}\\
 &\ge\varphi\!\left(\kappa_M\eta(z_*)\right)
 \notag\\
 &\ge
 \begin{cases}
  \varphi\!\left(
   \kappa_{\rm res}\beta_M\eta(z_*)
  \right),\\
  \varphi\!\left(
   d_{\rm missing}\eta(z_*)
  \right).
 \end{cases}
 \label{eq:tpc62-distinguished-lattice}
\end{align}
The number \(\kappa_M(z_*)\) is an exact Rayleigh quotient for this one
vector.  A favorable value of it cannot replace \(\kappa_M\) in
\cref{thm:tpc62-uniform-certificate-lattice}.
\end{proposition}

\begin{proof}
The first inequality is
\eqref{eq:tpc62-pointwise-q-certificate}; the equality is
\eqref{eq:tpc62-pointwise-factorization}.  Since
\(\kappa_M(z_*)\le\kappa_M\le
\kappa_{\rm res}\beta_M\) and
\(\kappa_M\le d_{\rm missing}\), monotonicity of \(\varphi\) gives all
remaining inequalities.  A single Rayleigh quotient supplies no Loewner
bound on other directions, so it cannot imply the uniform replacement.
\end{proof}

This distinction can be arbitrarily large.  Let \(A=I\),
\(\mathbf D=I\), and
\(M=\operatorname{diag}(0,L)\) on \(\C^2\).  At
\(z_*=e_1\), \(\kappa_M(z_*)=0\), while the uniform constant is
\(\kappa_M=L^2\).  Hence a perfect distinguished certificate may coexist
with an arbitrarily poor uniform one.

\subsection{How to read the lattice}

The exact shorted constant is the decision quantity for coefficient-
uniform reassembly.  The remaining nodes are sufficient lower
certificates.  Under the standing hypothesis
\(\eta_{\rm unif}<\infty\) of
\cref{thm:tpc62-uniform-certificate-lattice}, the partial order is
\begin{equation}
 \delta_{\rm sh}(A,M)
 \ge\varphi(q_{\rm unif})
 \ge\varphi(\kappa_M\eta_{\rm unif})
 \ge
 \begin{cases}
  \varphi(\kappa_{\rm res}\beta_M\eta_{\rm unif}),\\
  \varphi(d_{\rm missing}\eta_{\rm unif}).
 \end{cases}
 \label{eq:tpc62-certificate-poset-display}
\end{equation}
The two bottom nodes are not ordered.  Improving either one is useful only
as a way to bound \(\kappa_M\); once \(\kappa_M\) itself is known, neither
coarser number is charged in addition.  Likewise, a pointwise value
\(q(z_*)\) or \(\kappa_M(z_*)\) belongs only to the distinguished chain,
whereas \(q_{\rm unif}\), \(\kappa_M\), and
\(\delta_{\rm sh}\) belong to the uniform chain.

All conclusions here are L0 finite-dimensional identities and
inequalities.  They neither estimate the physical TPC constants nor show
that any of them is subpower in \(X\).  In particular, the lattice does
not establish cancellation in the affine M\"obius kernel, a fixed-shift
parity gain, a prime-pair lower square, or a twin-prime statement.
